\documentclass[lettersize,journal]{IEEEtran}
\usepackage{amsmath,amsfonts,amssymb}
\usepackage{array}
\usepackage{url}
\usepackage{booktabs}
\renewcommand{\thetable}{S\arabic{table}}
\renewcommand{\thesection}{S\arabic{section}}
\hyphenation{op-tical net-works semi-conduc-tor IEEE-Xplore}

\begin{document}

\title{Supplementary Material for\\ ``Certifying the Outcome of Autonomous
Vehicle Fallback Triggers: Closed-Loop Evaluation Under Latency and Dataset Shift''}
\author{Hassan~Waqas~Saleem, Rahat~Zeeshan~Saleem, and Hina~Maryam}
\markboth{Supplementary Material}{Saleem \MakeLowercase{\textit{et al.}}:
Certifying Fallback Triggers in Closed Loop}
\maketitle

This document contains the exploratory and post hoc analyses summarized
in Section~V of the main paper. All numbers use the same settings as the
main paper ($\alpha = 0.05$, $\delta = 0.10$, 1{,}994 Argoverse~2
calibration and 1{,}996 test scenarios, 1{,}504 nuScenes scenarios) and
can be regenerated from the released per-scenario records without
re-simulation.

\section{Class-Conditional Certification (Post Hoc)}
\label{sec:s-luo}
Table~\ref{tab:luo} lists the differences between the closest prior
work and this study.

\begin{table}[h]
\caption{Luo~\emph{et al.} versus this study.}
\label{tab:luo}
\centering
\scriptsize
\setlength{\tabcolsep}{3pt}
\begin{tabular}{p{1.9cm}p{2.85cm}p{2.85cm}}
\toprule
 & Luo~\emph{et al.} & This paper \\
\midrule
Guarantee & $\Pr[\text{no alert}\mid\text{unsafe}]\le\epsilon$, class-conditional & $\Pr[\text{harm not pre-empted}]\le\alpha$ over scenarios, high probability \\
Unsafe label & safety score of logged future $<f_0$ & nuPlan harm in the no-fallback run \\
Predictor & Trajectron++ & AutoBot \\
Alert consequence & none (scored on logs) & MRM executed in closed loop \\
Cost of alerting & false-positive rate & unnecessary stops, induced collisions, residual harm \\
Data & nuScenes, Lyft & Argoverse~2, nuScenes, Waymo \\
Shift & label shift only & dataset shift, weighted repair \\
Latency & not studied & studied \\
\bottomrule
\end{tabular}
\end{table}

Table~\ref{tab:luoresult} implements the certificate of Luo~\emph{et al.}
(cited in the main paper) in our framework: Mondrian split conformal
prediction on the unsafe class, where unsafe means harm in the reference
run and an alert means firing before it, with the finite-sample
correction and without their tie randomization. At $\epsilon=0.05$ on the
316 unsafe calibration scenarios it selects a threshold equivalent to
always firing, so the class-conditional target is vacuous for this
predictor at this level. The score separates harmful from harmless
scenarios well enough for a marginal target (AUROC 0.83) but not for a
5\,\% miss rate among harmful ones. At the marginal-matched
$\epsilon=\alpha/\pi=0.32$ it selects the same threshold as empirical
tuning. Deployed on nuScenes, that certificate has a conditional miss
rate of 66\,\% against its 32\,\% target. The harm base rate rises from
15.8\,\% to 22.1\,\%, but the conditional miss rate is far above what that
alone explains, so the class-conditional guarantee does not survive the
dataset shift either.

\begin{table}[h]
\caption{Marginal certificate (LTT) versus class-conditional certificates
(Luo-style), thresholds certified on Argoverse~2. Miss: marginal miss
rate; Cond.: miss rate among scenarios with harm in the reference run.
Rates in \%.}
\label{tab:luoresult}
\centering
\scriptsize
\setlength{\tabcolsep}{3pt}
\begin{tabular}{lcccccc}
\toprule
 & \multicolumn{3}{c}{Argoverse~2 test} & \multicolumn{3}{c}{nuScenes} \\
\cmidrule(lr){2-4}\cmidrule(lr){5-7}
Certificate & Miss & Cond. & Unnec. & Miss & Cond. & Unnec. \\
\midrule
LTT, marginal $\alpha=0.05$ & 3.7 & 24.2 & 28.5 & 13.2 & 58.1 & 26.1 \\
Luo-style, $\epsilon=0.05$ & 0.0 & 0.0 & 84.9 & 5.0 & 19.6 & 77.9 \\
Luo-style, $\epsilon=0.32$ & 4.3 & 28.5 & 26.0 & 14.9 & 66.0 & 21.3 \\
\bottomrule
\end{tabular}
\end{table}

\section{Reverse Shift (Exploratory, Post Hoc)}
\label{sec:s-reverse}
We also certified on half of the nuScenes scenarios and deployed on
Argoverse~2 test, a direction pre-registered as a reverse-shift analysis.
As pre-registered, with all scenarios, the procedure \emph{refused to
certify} on both halves, because of the 5.0\,\% floor. On the avoidable
subset (post hoc), thresholds certified on nuScenes ($n\!\approx\!715$)
give an Argoverse~2 miss rate of 1.9\,\% and 2.0\,\%, both below the
target, at 40.6\,\% and 39.5\,\% unnecessary stops. A certificate from the
harder dataset therefore transfers safely but conservatively, while the
reverse does not transfer. This asymmetry is a single-pair observation
pending replication.

\section{Recalibration With a Few Labeled Target Scenes}
\label{sec:s-fewshot}
Table~\ref{tab:recal} draws $k$ labeled target scenes at random and
evaluates on the rest. The nuScenes analysis (avoidable subset, 60 draws)
was exploratory and not pre-registered. Its Waymo counterpart (all
scenarios, 200 draws) was pre-registered as R1--R4 before any Waymo
outcome was read, and its outcomes are reported in Section~V of the main
paper. On nuScenes, uncertified tuning on 40--60 target scenes exceeds the
target in about two draws in three. A certified threshold on the same
scenes keeps the target (violations 0--10\,\%) but stops in 53--81\,\% of
scenarios, against 27\,\% without target labels. A handful of labeled
scenes therefore restores a \emph{valid} certificate, not the original
operating point. On Waymo, uncertified tuning violates the target in
69.5\,\% and 59\,\% of draws with 40 and 60 target scenes, and weighting
the Argoverse~2 calibration set with 60 Waymo scenes
($n_{\mathrm{eff}} = 232$) keeps validity at 0\,\% violation.

\begin{table}[h]
\caption{Recalibration with $k = 60$ labeled target scenes. nuScenes:
exploratory, avoidable subset, 60 draws. Waymo: pre-registered
(R1--R4), all scenarios, 200 draws. Viol.: fraction of draws with test
miss $>\alpha$. Rates in \%.}
\label{tab:recal}
\centering
\scriptsize\setlength{\tabcolsep}{2.5pt}
\begin{tabular}{lcccccc}
\toprule
& \multicolumn{3}{c}{nuScenes} & \multicolumn{3}{c}{Waymo} \\
\cmidrule(lr){2-4}\cmidrule(lr){5-7}
Method & Miss & Viol. & Unnec. & Miss & Viol. & Unnec. \\
\midrule
AV2-certified, no labels & 9.2 & -- & 27.1 & 2.7 & 0 & 45.0 \\
Target-only tuned & 6.0 & 70 & 36.7 & 6.2 & 59 & 29.6 \\
Target-only LTT & 0.8 & 0 & 72.5 & 1.5 & 2.5 & 58.6 \\
AV2 + 60, weighted LTT & 2.8 & 0 & 53.3 & 1.7 & 0 & 49.8 \\
\bottomrule
\end{tabular}
\end{table}

\section{Simulator Sensitivity}
\label{sec:s-sens}
Three changes to the simulator, each paired by scenario seed against the
main run (reactive traffic, 4.0\,m/s$^2$, 2\,Hz), leave every conclusion
of the main paper unchanged. Non-reactive log-replay traffic raises the
harm base rate (17.7\,\% versus the main run's 15.1\,\%, since logged
agents cannot react to an ego that deviates from its log) and, with it,
unnecessary stops by 7.9 points [95\,\% CI 5.3, 10.8]. The miss rate
itself does not move (paired difference $-0.5$ points [$-1.3$, 0.0]). A
2.0\,m/s$^2$ comfort deceleration changes neither miss nor stops (paired
difference exactly 0 in both), because both alert-level quantities depend
only on the reference run and the firing tick (Observation~1 of the main
paper), not on how hard the vehicle then brakes. It does lower induced
collisions (2.6\,\% to 1.6\,\%), the one quantity that does depend on
post-fire kinematics. Decisions at 10\,Hz on 100 scenarios give a miss
rate of 2.9\,\% [1.0, 8.2], consistent with the main run within an
interval widened by the smaller sample.

\section{Certifying Residual Harm Directly (Post Hoc)}
\label{sec:s-outcome}
\emph{Update.} The attempt described here, with a threshold ordered by the monotone alert-level order and a 5\,\% target, failed. A later amendment (Amendment~11) certified the same quantity at a 7.5\,\% target with thresholds ordered by out-of-fold estimated risk, and succeeded (Section~V-D of the main paper). The failure below is therefore attributable to the target and the order, not to outcome-level certification as such.

The certificate in the main paper bounds the alert indicator
$L(X,\lambda)$, not harm in the run the vehicle then drives. We therefore
tried to certify \emph{residual harm} $H(X,\lambda)$, the indicator that
the nuPlan harm definition fires anywhere in the executed run, before or
after the fallback engages. $H$ is not monotone in $\lambda$, because
firing earlier can itself create harm (4.5\,\% of fresh calibration
scenarios have a non-monotone $H$), so Conformal Risk Control does not
apply by construction. Learn-then-Test needs no monotonicity but is
tested on the same fresh Argoverse~2 split as the replication analysis of
the main paper ($n_{\mathrm{cal}} = n_{\mathrm{test}} = 904$, target
$\alpha_H = 0.05$, $\delta = 0.10$). The comparisons were pre-registered
as N1--N5 before running.

\begin{table}[h]
\caption{Certifying residual harm on the fresh split. Rates in \%, Wilson
95\,\% intervals. Viol.: fraction of 200 resplits whose test residual
harm exceeds 5\,\%. The second predictor gives the same picture
(outcome-level LTT: 2.8\,\% harm, 84.8\,\% stops, 16.0\,\% violations).}
\label{tab:outcome}
\centering
\scriptsize\setlength{\tabcolsep}{2.5pt}
\begin{tabular}{lcccc}
\toprule
Certificate on & Resid. harm [CI] & Miss & Stops & Viol. \\
\midrule
Alert $L$, LTT (main paper) & 5.0 [3.7, 6.6] & 3.4 & 29.0 & -- \\
Residual harm $H$, LTT & 2.8 [1.9, 4.1] & -- & 84.8 & 15.0 \\
Residual harm $H$, CRC & 3.9 [2.8, 5.3] & -- & 37.5 & 45.5 \\
\bottomrule
\end{tabular}
\end{table}

Three results, each contrary to what we pre-registered, follow from
Table~\ref{tab:outcome}. First, certifying $H$ at 5\,\% with LTT selects a
threshold at which the vehicle fires in essentially every scenario (84.8\,\%
unnecessary stops), and even then it exceeds the target in 15\,\% of
resplits, above $\delta$, so it is not a valid certificate (N1 refuted).
Second, the alert-level threshold's residual harm on this fresh split is
5.0\,\% [3.7, 6.6], so its interval includes the 5\,\% target, unlike the
6.5\,\% on the main split (N2 refuted on fresh data; the interval is
wide enough that a modest excess is not excluded). Third, CRC violates in
45.5\,\% of resplits, but LTT also exceeds $\delta$, so a
non-monotone-aware calibrator is not by itself sufficient (N3 refuted). A
joint certificate on residual harm and induced collisions over threshold
and two MRM decelerations (2 and 4\,m/s$^2$) found no configuration other
than always firing (N4 refuted). The practical reading is that, with this
score and this maneuver, residual harm of 5\,\% or less is not reachable
by tuning the trigger threshold alone at a useful operating point. A
different maneuver or score is needed, not a different calibrator.

\section{Shift Severity and the Cost of the Intervention (Descriptive)}
\label{sec:s-cost}
These two tables are computed from the stored records and test no
hypothesis.

\emph{Shift severity.} Table~\ref{tab:severity} measures how far each
target is from the Argoverse~2 calibration set in the seven
pre-deployment covariates used for the density ratio. The domain
classifier is evaluated by 5-fold cross-validated AUROC (0.5 means
indistinguishable), and the Argoverse~2 test split is a null control.
nuScenes is much easier to tell apart than Waymo, has a Kish effective
sample size of 66 against 288, and a 5.0\,\% floor of harm that no trigger
can pre-empt. The certificate failed on the more distant target and held
on the nearer one, which is consistent with shift severity mattering,
but two targets cannot establish that, and the floor alone accounts for
5.0 of the 13.2 points missed on nuScenes.

\begin{table}[h]
\caption{Shift severity and the Argoverse~2-certified threshold
($\lambda = 1.47$). AUROC: cross-validated domain classifier against
the Argoverse~2 calibration set. Floor: harm present even when the
fallback fires at the first decision tick. Rates in \%.}
\label{tab:severity}
\centering
\scriptsize\setlength{\tabcolsep}{2.5pt}
\begin{tabular}{lcccccc}
\toprule
Target & $n$ & AUROC & $n_{\mathrm{eff}}$ & Harm & Floor & Miss [CI] \\
\midrule
AV2 test (control) & 1{,}996 & 0.50 & 1{,}978 & 15.1 & 0.0 & 3.7 [2.9, 4.6] \\
Waymo & 1{,}000 & 0.91 & 288 & 15.4 & 0.1 & 2.7 [1.9, 3.9] \\
nuScenes & 1{,}504 & 0.97 & 66 & 22.1 & 5.0 & 13.2 [11.6, 15.0] \\
\bottomrule
\end{tabular}
\end{table}

\emph{Scenarios from the same scene are not independent.} The Wilson
intervals in the main paper treat scenarios as independent. The nuScenes
scenario identifiers contain the scene name, which allows a check:
the 1{,}504 nuScenes scenarios in the main analysis come from only 37
scenes (41 per scene on average). Resampling whole scenes
(cluster bootstrap, 10{,}000 draws) widens the interval for the miss rate
of the Argoverse~2-certified threshold from [11.6, 15.0] to [6.1, 21.0],
which still excludes the 5\,\% target, but only just. A second, largely
disjoint set of nuScenes scenarios (1{,}342 scenarios from 23 scenes, one
scene in common, collected for the replication analyses) gives a miss rate
of 9.5\,\% with Wilson interval [8.1, 11.2] and cluster-bootstrap interval
[3.9, 16.5], which includes the target. The first-tick floor also differs
between the two samples (5.0\,\% against 0.0\,\%), so it varies strongly
with which scenes are drawn. The direction of the failure therefore
replicates, but with a few dozen scenes per sample the size of the
nuScenes failure is far less certain than the main-paper intervals
suggest. Argoverse~2 and Waymo identifiers are opaque, so the same
check, and group-level resplits for the validity analysis, were not
possible there.

\emph{Where the induced collisions come from.} At the certified threshold
on Argoverse~2 test, 2.6\,\% of scenarios [1.9, 3.3] suffer an induced
collision. All 129 induced contact events are other vehicles striking the
ego, 99 from behind and 30 from the side. They concentrate where
the ego is slow or fires early. Scenarios in the two lower terciles of
initial ego speed (the lowest starts at standstill, the middle one below
4.6\,m/s) have induced rates of 3.6\,\% [2.4, 5.3]
and 3.5\,\% [2.3, 5.1], against 0.6\,\% [0.2, 1.5] in the upper tercile, and
scenarios in which the fallback fires at or before the median firing tick
(tick 5, 2.5\,s) have 9.2\,\% [6.8, 12.4] against 3.0\,\% [1.7, 5.1]
for later firing. These contacts do not look like a follower caught off
guard by hard braking. Of the 129, only 3.9\,\% occur within 2\,s of the
fallback firing, the median is 8.3\,s after it, and 78\,\% come more than
5\,s after it, by which time a vehicle braking at 4\,m/s$^2$ from city speed
has long been stationary. The 129 contacts come from only 51 scenarios,
and 67\,\% of them are in scenarios where the ego starts at standstill, so
the fallback there changes little about what the ego does. We therefore
cannot attribute this cost to the act of stopping. It may instead come
from simulated traffic, driven along logged paths with a car-following
model, that does not yield to a vehicle which stays stationary in its path
for many seconds, which a real driver would also be expected to avoid
(released with the code). The 2.6\,\%
induced rate should be read as an upper bound on the cost that stopping
itself causes in this simulator.

To grade severity we re-ran the single forced rollout of each of the 51
scenarios with instrumentation (Amendment~9, descriptive). The re-run
reproduces the stored induced-contact steps exactly in 34 scenarios and
the first induced contact in 49, but later contacts differ in 17, so
the contact dynamics after a first impact are not run-to-run reproducible,
which is one more way in which the forced-rollout table is exact only up
to the first contact. At the first induced contact the ego's speed has a
median of 0.6\,m/s (27\,\% below 0.1\,m/s and 65\,\% below 1\,m/s), the other
vehicle's speed a median of 3.6\,m/s, and the speed difference a median of
2.4\,m/s (quartiles 1.2 and 5.8, maximum 11.8; this is a scalar speed
difference, not a closing speed along the contact axis). Nineteen of the
51 first contacts are lateral and 32 are from behind, a median 6.5\,s after
firing. At the firing step an agent was in the ego's lane behind it in 39
of the 51 scenarios, at a median gap of 8.3\,m and a median speed of
0\,m/s. These are therefore low- to moderate-speed contacts with a slow or
stationary ego by traffic that keeps moving along its logged path, not
hard-braking surprises. The fallback also costs progress: where it fires,
the vehicle stops for 7.4\,s on average and completes 49\,\% of the route,
against 94\,\% where it does not. The rates count scenarios, not contacts.

\section{Direct Execution Against the Lookup}
\label{sec:s-direct}
The main paper evaluates every trigger by lookup in a table of forced-fire
rollouts, which is exact when a run under the trigger coincides with the
reference run until it fires (Observation~1). The assumption is
measured in the forced runs but was not tested by executing a trigger.
This check, pre-registered as Amendment~7 and extended in Amendment~9,
does so. It draws two independent random samples of stored Argoverse~2
test scenarios (120 and 140 scenarios, 252 distinct after merging by
seed) and runs four geometric triggers live, with the predictor scoring at
every decision tick and no score replay or forced-fire table: the
certified LTT threshold with zero and with one decision tick of latency,
and the thresholds at the 25th and 75th percentile of the calibration peak
scores (about 75\,\% and 25\,\% of scenarios fire). Each run is scored with
the same nuPlan harm definition and compared, scenario by scenario, with
the lookup at the same threshold and latency. No direct run failed.
(The first draw was first analyzed as a snapshot of 61 scenarios, because
the stop signal described in Amendment~7 did not take effect and the job
ran to completion. The figures here supersede that snapshot.)

\begin{table}[h]
\caption{Direct execution versus lookup, 252 scenarios $\times$ 4
triggers = 1{,}008 pairs. Disagreement: pairs in which the two differ.}
\label{tab:direct}
\centering
\scriptsize\setlength{\tabcolsep}{3pt}
\begin{tabular}{lcc}
\toprule
Quantity & Disagree & Wilson 95\,\% upper bound \\
\midrule
Firing step (or no firing) & 0 / 1{,}008 & 0.4\,\% \\
Missed intervention & 0 / 1{,}008 & 0.4\,\% \\
Unnecessary stop & 0 / 1{,}008 & 0.4\,\% \\
Induced collision & 0 / 1{,}008 & 0.4\,\% \\
Harm anywhere in the run & 1 / 1{,}008 & 0.6\,\% \\
\bottomrule
\end{tabular}
\end{table}

Table~\ref{tab:direct} shows no disagreement in any alert-level quantity
or in the induced collisions. The aggregate rates are identical for every
trigger (for example 3.6\,\% missed and 27.8\,\% unnecessary stops for the
certified threshold, and 4.8\,\% missed with one tick of latency). One pair
differs in harm anywhere in the run (seed 1934, certified threshold): the
lookup records harm after the fallback and the direct run does not. We did
not trace its cause. The re-run of induced collisions in
Section~S6 shows that contact events after a first impact are not always
reproducible from run to run, which is a plausible source. Route
completion is identical in 454 of the 1{,}008 pairs and differs by 0.0002
on average and 0.036 at most. The pre-registered rule, that the lookup is
refuted as a substitute if any quantity disagrees in more than 5\,\% of
pairs, is therefore not triggered, and the upper bound on the disagreement
rate of the alert-level quantities is 0.4\,\%. The check covers one
dataset, one predictor and one maneuver, and does not test stochastic
traffic, for which assumption (A1) would not hold and the lookup would
estimate the outcome under a different random draw than the one the
trigger would have met. The magnitudes of the pre-fire deviations in the
forced runs were not stored, only whether they exceeded 1\,cm.

\section{Replication on Fresh Data and a Second Predictor}
\label{sec:s-replicate}
Table~\ref{tab:replicate} gives the numbers behind the replication analysis in the main paper.
Violation frequencies over 200 resplits have Wilson 95\,\% intervals of [4.6, 12.0] for LTT with
the primary predictor and [5.8, 13.8] with the second, so, as in the main split, the point
estimates meet $\delta = 0.10$ but their Monte Carlo intervals do not exclude a violation
frequency above it. The uncertified baselines (46--51.5\,\%) are far outside it.
The gap between the geometric and the confidence trigger also replicates:
with the same target, LTT stops 29.0\,\% [26.1, 32.0] of scenarios against
67.0\,\% [63.9, 70.0] for the tuned confidence trigger on the fresh split,
and 30.6\,\% [27.7, 33.7] against 68.3\,\% [65.1, 71.2] with the second
predictor.

\begin{table}[h]
\caption{Replication on fresh, disjoint data ($n_{\mathrm{cal}} =
n_{\mathrm{test}} = 904$). Violation: fraction of 200 resplits with test
miss $>\alpha$.}
\label{tab:replicate}
\centering
\scriptsize\setlength{\tabcolsep}{3pt}
\begin{tabular}{lcccc}
\toprule
& \multicolumn{2}{c}{Fresh held-out} & \multicolumn{2}{c}{2nd predictor} \\
\cmidrule(lr){2-3}\cmidrule(lr){4-5}
Method & Miss (\%) & Viol. (\%) & Miss (\%) & Viol. (\%) \\
\midrule
LTT (geometric) & 3.4 & 7.5 & 3.9 & 9.0 \\
CRC (geometric) & 4.3 & 46.0 & 4.6 & 47.5 \\
Tuned geometric & 4.3 & 51.5 & 4.6 & 51.0 \\
T1 confidence & 4.8 & 50.0 & 4.6 & 46.0 \\
\bottomrule
\end{tabular}
\end{table}

\emph{Ensemble baseline with the second predictor.} The replication above
lacked the Monte Carlo dropout trigger. We re-ran only the reference run
of every fresh calibration and test scenario with five dropout passes
(Amendment~9), reusing the forced-fire outcomes, which do not depend on
any score. The re-run reproduced the stored reference (tick count and
geometric trace) in 865 of 904 calibration and 867 of 904 test scenarios,
and we dropped the others (4.3\,\% and 4.1\,\%), so these figures use 865
and 867 scenarios. On the test scenarios the certified geometric trigger
misses 3.2\,\% [2.2, 4.6] with 26.4\,\% unnecessary stops, the tuned
confidence trigger 5.0\,\% [3.7, 6.6] with 66.2\,\%, and the tuned ensemble
trigger 5.7\,\% [4.3, 7.4] with 59.6\,\%. The tuned ensemble trigger
therefore misses the 5\,\% target, which the pre-registered comparison
``among triggers that meet the target'' did not anticipate, so the rule
as written cannot be evaluated for it. Unconditionally, the certified
trigger stops 33.2 points less than the ensemble trigger (paired 95\,\%
CI $[-37.5, -29.0]$) and 39.8 points less than the confidence trigger
($[-43.7, -35.8]$), with a lower miss rate than both. The
geometric-versus-ensemble gap is thus reproduced with a second predictor,
and the open-loop conformal comparison alone remains untested there.


\section{Certifying Under the Deployed Latency}
\label{sec:s-latency}
The main paper shows that a certificate calibrated at zero latency fails
when the actuation chain is slower. Because the loss at latency $\ell$
decision ticks is the stored miss indicator at tick $\tau(X)+\ell$, LTT can
be run directly on that loss without new simulation (Amendment~8,
pre-registered before the run). We certify on the Argoverse~2 calibration
set at $\ell \in \{0,1,2\}$ ticks (0, 0.5, 1.0\,s), deploy at the same
latency on the test set, and repeat over the same 200 half-and-half
resplits as the main paper.

\begin{table}[h]
\caption{LTT calibrated at the deployed latency (``aware'') versus the
zero-latency certificate deployed at latency $\ell$ (``naive''),
Argoverse~2 test. Miss and stops on the test set, rates in \%.
Viol.: fraction of 200 resplits with test miss above 5\,\%, with Wilson
95\,\% intervals.}
\label{tab:latency}
\centering
\scriptsize\setlength{\tabcolsep}{2.5pt}
\begin{tabular}{cccccc}
\toprule
$\ell$ & Naive miss & Naive viol. & Aware miss & Aware stops & Aware viol. \\
\midrule
0 & 3.7 & 9.5 [6.2, 14.4] & 3.7 & 28.5 & 9.5 [6.2, 14.4] \\
1 & 5.3 & 92.0 [87.4, 95.0] & 3.6 & 34.4 & 12.0 [8.2, 17.2] \\
2 & 9.1 & 100 [98.1, 100] & 4.1 & 42.5 & 11.5 [7.8, 16.7] \\
\bottomrule
\end{tabular}
\end{table}

The zero-latency certificate is not valid under latency, as claimed.
The latency-aware certificate brings the miss rate back below 5\,\% on the
test set, at a cost of 6 and 14 points of unnecessary stops. Its
violation frequency, however, is 12.0\,\% and 11.5\,\% at one and two ticks,
above the $\delta = 0.10$ we pre-registered, so by our own rule the claim
that latency-aware certification restores validity was refuted under the uncorrected count. Under the corrected count of Amendment~11 (Section~S17), 1.5\,\% and 0.5\,\%, the certificate is indeterminate, not violated, and the pre-registered re-reading holds. Three
readings are possible and we cannot separate them. The Wilson intervals
include 0.10. The statistic counts resplits whose finite test half
exceeds 5\,\%, not resplits whose population risk does, and the
zero-latency case is already 9.5\,\%. And the loss at a positive latency
is less smooth in $\lambda$, which may weaken the fixed-sequence test in
practice. The guarantee itself is unchanged, since the loss is still a
fixed function of the scenario and the threshold, but we report the
frequency as measured. Across resplits the mean unnecessary-stop rate of
the latency-aware certificate is 27.6, 31.8 and 57.0\,\%, so at two ticks
the cost depends strongly on the calibration draw.

\section{Natural Shifts Inside Argoverse 2: Leave-One-City-Out}
\label{sec:s-city}
The main paper's shift evidence rests on two target datasets. Argoverse~2
labels each scenario with one of six cities, which gives six further
natural shifts at no simulation cost. For each city we run LTT
($\alpha = 0.05$, $\delta = 0.10$, geometric score) on all stored
calibration and test scenarios outside that city (about 3{,}000 to
3{,}800 scenarios) and deploy the threshold on the held-out city. We
pre-registered (Amendment~9) that the certificate would count as
transferring if the held-out miss rate were at most 5\,\% in at least 80\,\%
of the cities.

\begin{table}[h]
\caption{Leave-one-city-out on Argoverse~2 (pooled calibration and test
scenarios, $n = 3{,}990$). Harm: harm rate of the held-out city. Miss and
stops: held-out city at the threshold certified on the other cities, with
Wilson 95\,\% intervals for the miss rate. Tuned: miss of the uncertified
threshold tuned on the other cities. Rates in \%.}
\label{tab:city}
\centering
\scriptsize\setlength{\tabcolsep}{2.5pt}
\begin{tabular}{lccccc}
\toprule
Held-out city & $n$ & Harm & LTT miss [CI] & LTT stops & Tuned miss \\
\midrule
Dearborn & 546 & 8.4 & 2.0 [1.1, 3.6] & 17.0 & 2.9 \\
Palo Alto & 229 & 12.7 & 4.4 [2.4, 7.9] & 21.0 & 4.8 \\
Austin & 931 & 14.2 & 3.7 [2.6, 5.1] & 26.7 & 4.1 \\
Pittsburgh & 870 & 15.4 & 4.4 [3.2, 5.9] & 30.7 & 5.1 \\
Washington, DC & 450 & 19.3 & 5.3 [3.6, 7.8] & 26.7 & 6.2 \\
Miami & 964 & 19.7 & 6.1 [4.8, 7.8] & 30.2 & 7.0 \\
\bottomrule
\end{tabular}
\end{table}

The certificate met the target in four of the six cities (67\,\%), below
the 80\,\% we required, so the pre-registered claim is refuted. The
uncertified threshold met it in three. The two failures, Miami and
Washington, DC, are the two cities with the highest harm rates, and in
only Miami does the interval exclude 5\,\%. The failure is not an artifact
of holding the city out: the threshold certified on all cities, which
saw them, also misses 5.2\,\% and 5.1\,\% there. This is what a
\emph{marginal} guarantee implies, since the miss rate is the harm rate
times the fraction of harmful scenarios left unalerted. Post hoc and
descriptive, across these six cities and the two datasets of the main
paper, the held-out miss rate and the harm rate of the target have a
Spearman correlation of 0.90 ($n = 8$, so this is suggestive only). A
shift in the base rate of harm, which an unlabeled-covariate reweighting
cannot see, is therefore a plausible common cause of the failures, and
it would also explain why the certificate failed on nuScenes (22.1\,\%
harm) but held on Waymo (15.4\,\%).

\section{Shift in the Traffic Behavior Model}
\label{sec:s-behavior}
All main results use reactive IDM traffic. To test a different behavior
model we deploy the thresholds certified on reactive calibration data on
the non-reactive log-replay arm, paired by scenario seed with the reactive
test arm (378 scenarios with complete records). The harm base rate is
16.7\,\% under reactive traffic and 17.7\,\% under replay. The certified
geometric trigger misses 3.7\,\% [2.2, 6.1] on replay traffic, so the
pre-registered claim that it still meets the 5\,\% target is not
refuted. The tuned confidence trigger also meets the target (3.4\,\%), but at
71.7\,\% unnecessary stops against 35.4\,\% for the certified geometric
trigger (paired difference $-36.2$ points, 95\,\% CI $[-41.8, -30.7]$). The
stop-rate advantage therefore persists under a second traffic model, although
both are rule-based and neither is a learned agent model.

\section{Spatial-Group Dependence}
\label{sec:s-spatial}
Scenarios from the same place or drive may be dependent, which would
weaken the exchangeability that the guarantee needs. The converted data
carry no recording identifier for Argoverse~2 and Waymo, so
(Amendment~10, pre-registered) we group scenarios by the ego's initial
position: single-linkage components of scenarios whose ego start points lie
within 150\,m of each other, formed separately for each Argoverse~2 city.
This is a proxy, not a recording identifier. It gives 547 groups over the
3{,}990 pooled Argoverse~2 scenarios (the largest has 649, 16\,\% of them, the
median group is a single scenario) and 664 groups over the 1{,}000 Waymo
scenarios (the largest has 13).

We repeat the main validity analysis with whole groups assigned to
calibration or test (200 resplits, calibration share about 53\,\%). The
violation frequency of LTT is 17.0\,\% (Wilson 95\,\% interval [12.4, 22.8])
against 9.5\,\% for random resplits, and the mean test miss rate is 4.2\,\%.
The pre-registered claim that it stays at or below $\delta = 0.10$ is
therefore refuted under the uncorrected count (it is indeterminate, at 1.0\,\%, under the corrected count of Section~S17). A post hoc sensitivity at 50\,m, which gives 1{,}796
groups with the largest at 79 scenarios, yields 14.0\,\% [9.9, 19.5].
Dependence between nearby scenarios is thus one more reason the
high-probability statement is weaker in practice than in theory, although
random and group-wise splits differ by less than the Monte Carlo error of
either would suggest at 200 resplits.

The miss-rate intervals are far less sensitive. A group bootstrap of the
certified threshold's miss rate gives [3.0, 4.5]\,\% on the Argoverse~2
test scenarios (360 groups; Wilson [2.9, 4.6]) and [1.7, 3.8]\,\% on Waymo
(664 groups; Wilson [1.9, 3.9]). Place-level dependence therefore does not
change the point conclusions about the miss rate, including the
conservative result on Waymo.

\section{A Second Architecture: Wayformer}
\label{sec:s-wayformer}
Every main result uses AutoBot, and the second predictor of
Section~S8 is also an AutoBot. To test whether the findings are specific
to the architecture, we repeated the full closed-loop campaign
(Amendment~10, pre-registered) on the same 904 calibration and 904 test
scenarios with a Wayformer forecaster from the same framework (minADE$_6 =
0.967$, a partly trained checkpoint, the same 2.1\,s history and 6\,s
horizon). A check on identical simulator states found a median 3\,s
endpoint disagreement of 1.6\,m with AutoBot, so the wrapper produces
comparable forecasts. T2 and T3 are not repeated.

\begin{table}[h]
\caption{Wayformer arm, Argoverse~2 fresh test scenarios ($n = 904$). Viol.:
fraction of 200 resplits with test miss above 5\,\%. Rates in \%.}
\label{tab:way}
\centering
\scriptsize\setlength{\tabcolsep}{3pt}
\begin{tabular}{lccc}
\toprule
Method & Miss [95\,\% CI] & Unnec. stops & Viol. \\
\midrule
LTT (geometric) & 3.1 [2.2, 4.4] & 41.8 & 7.0 \\
Tuned geometric & 3.8 [2.7, 5.2] & 36.7 & 48.0 \\
CRC (geometric) & -- & -- & 44.0 \\
Tuned T1 confidence & 4.2 [3.1, 5.7] & 67.8 & 53.0 \\
\bottomrule
\end{tabular}
\end{table}

Both parts of the pre-registered claim hold. LTT is valid (7.0\,\%
violations, Wilson interval [4.2, 11.4], and a test miss of 3.1\,\%), and
the certified geometric trigger stops 26.0 points less than the tuned
confidence trigger (paired 95\,\% CI $[-29.8, -22.1]$), which meets the
target with 4.2\,\% missed. The uncertified geometric threshold and CRC
violate the target in 48\,\% and 44\,\% of resplits, as with AutoBot. The
certificate itself and the geometric advantage therefore do not depend on
the architecture. The price does depend on the predictor: with Wayformer
the certified trigger stops 41.8\,\% of scenarios, against 26 to 29\,\% with
the two AutoBot models, which is consistent with a less accurate forecaster
needing a more conservative threshold to meet the same target.

\section{Physics Baselines, Ablations and Open-Loop Ranking}
\label{sec:s-physics}
The main paper's Table~I adds a time-to-collision (TTC) trigger, certified with the same procedure. Because the harm label is itself defined by time-to-collision, the TTC trigger is close to the label, so we repeat the comparison under the alternative harm definitions of Section~V (post hoc, not pre-registered). The certified TTC trigger needs fewer unnecessary stops than the predictor-based geometric trigger under three of the four definitions: collision only ($\alpha = 0.01$) 18.8\,\% against 45.1\,\%, TTC $<0.5$\,s 1.6\,\% against 15.8\,\%, the headline TTC $<0.95$\,s 7.1\,\% against 28.5\,\%. Under TTC $<1.5$\,s the certified TTC threshold is very conservative (77.5\,\% against 34.0\,\%). A headway-distance trigger is worse than the geometric trigger under all four. The ablations below and the open-loop ranking behind Section~V-B follow.

\begin{table}[h]
\caption{Ablations on the Argoverse~2 test set, certified by LTT.}
\label{tab:abl}
\centering
\footnotesize
\begin{tabular}{lcc}
\toprule
Variant & Miss (\%) & Unnec. stops (\%) \\
\midrule
Full score, LTT, no latency & 3.7 & 28.5 \\
Route-following ego corridor & 0.0 & 84.9 \\
Width-only isotropic radius & 3.6 & 27.5 \\
Confidence instead of geometry & 3.6 & 75.8 \\
Ensemble spread instead of geometry & 3.6 & 72.3 \\
CRC instead of LTT & 4.3 & 26.0 \\
Empirical tuning instead of LTT & 4.3 & 26.0 \\
+1 decision tick latency (0.5\,s) & \textbf{5.3} & 26.3 \\
+2 decision ticks latency (1.0\,s) & \textbf{9.1} & 23.4 \\
\bottomrule
\end{tabular}
\end{table}

\begin{table}[h]
\caption{Open-loop AUROC versus closed-loop unnecessary stops at the
certified threshold, Argoverse~2 test.}
\label{tab:h1}
\centering
\footnotesize
\begin{tabular}{lcc}
\toprule
Score & AUROC & Unnec. stops (\%) \\
\midrule
Geometric, isotropic radius & 0.843 & 27.5 \\
Geometric, length-aware box & 0.832 & 28.5 \\
Predictor gap & 0.817 & 31.8 \\
Geometric, route corridor & 0.577 & 84.9 \\
Confidence & 0.448 & 75.8 \\
Ensemble spread & 0.442 & 72.3 \\
\bottomrule
\end{tabular}
\end{table}

\begin{table}[h]
\caption{LTT (geometric) versus the better of T1/T2 under alternative
harm definitions, Argoverse~2 test. Unnecessary-stop rates in \%;
``Best'' is the better of T1/T2 among those meeting the miss target, and
the paired 95\,\% CI is of the stop-rate difference.}
\label{tab:align}
\centering
\scriptsize\setlength{\tabcolsep}{2pt}
\begin{tabular}{lcccc}
\toprule
Harm definition & LTT & Best & Which & Diff. [95\,\% CI] \\
\midrule
Collision only, $\alpha{=}0.01$ & 45.1 & 83.9 & T2 & $-38.8$ [$-41.4,-36.2$] \\
TTC $<$ 0.5\,s & 15.8 & 38.6 & T1 & $-22.7$ [$-25.5,-20.1$] \\
TTC $<$ 0.95\,s (headline) & 28.5 & 68.3 & T1 & $-39.8$ [$-42.4,-37.2$] \\
TTC $<$ 1.5\,s & 34.0 & 68.3 & T2 & $-34.4$ [$-36.8,-31.9$] \\
\bottomrule
\end{tabular}
\end{table}

\section{Counterfactual Trigger: Safe Windows, Shift and Prevalence Correction}
\label{sec:s-cswc}
\emph{Safe windows.} On the 1{,}994 training scenarios, 15.8\,\% have harm in the
no-fallback run. Harm is unavoidable (empty safe window) in 0.35\,\% of
scenarios, so almost all harm is preventable by some firing tick. In
12.7\,\% of scenarios the window is not contiguous (a safe tick, an unsafe
tick, then a safe tick again), which is why the outcome loss is certified
directly instead of through a union bound on too-late and too-early
firing.

\emph{Prevalence-corrected certificate (refuted).} A marginal miss rate is the
harm rate times the miss rate among harmful scenarios, and the held-out
miss rate tracked the target's harm rate. We therefore estimated each
target's harm rate from unlabeled data with black-box shift estimation (the
pilot alert as classifier) and certified a class-conditional miss rate at
$\alpha$ over the estimate's upper bound (Amendment~11, N-A). The estimate
was unreliable (Table~\ref{tab:prev}): the alert rate shifts for reasons
other than harm rate, so the label-shift assumption fails. The corrected
certificate met the marginal target on 5 of 8 targets, against 6 of 8 for
plain LTT, and the pre-registered claim (at least 7 of 8 and more than
plain LTT) is refuted.

\begin{table}[h]
\caption{Black-box shift estimate of the target harm rate and marginal miss
at the target, unweighted LTT versus the prevalence-corrected certificate
(PC). Rates in \%.}
\label{tab:prev}
\centering
\scriptsize\setlength{\tabcolsep}{2.5pt}
\begin{tabular}{lccccc}
\toprule
Target & True harm & Estimate & Upper bound & LTT miss & PC miss \\
\midrule
Austin & 14.7 & 8.7 & 13.4 & 4.1 & 4.9 \\
Dearborn & 8.4 & 0.0 & 0.0 & 2.2 & 8.4 \\
Miami & 20.0 & 33.9 & 38.4 & 6.1 & 1.9 \\
Palo Alto & 11.6 & 0.0 & 6.5 & 3.9 & 7.7 \\
Pittsburgh & 15.5 & 25.1 & 29.5 & 4.6 & 2.0 \\
Washington, DC & 17.6 & 21.8 & 28.1 & 4.2 & 2.1 \\
nuScenes & 22.1 & 4.8 & 9.4 & 14.0 & 16.3 \\
Waymo & 15.4 & 46.3 & 51.7 & 2.8 & 0.5 \\
\bottomrule
\end{tabular}
\end{table}

\emph{Counterfactual trigger under shift (refuted).} Applying the same
correction to the outcome target of the learned trigger (tightening
7.5\,\% by the ratio of the source harm rate to the upper bound) made it
meet the target on 4 of 8 targets and refuse to certify on 3, against 6 of 8
without the correction (fails without it: Miami 10.8\,\%, nuScenes 8.2\,\%).
The pre-registered claim that the correction helps is refuted, for the
same reason.

\section{Outcome Audit With an Upper Confidence Bound}
\label{sec:s-audit}
The alert-level threshold (Learn-then-Test on the Argoverse~2 calibration
split, 5\,\% target) was applied to a fresh split of 904 scenarios, where
its residual harm is 6.7\,\% and total harm (with induced collisions) 8.1\,\%.
One-sided 90\,\% Clopper--Pearson upper bounds are 7.9\,\% and 9.4\,\%. The
realized harm on each target is compared with the bound in
Table~\ref{tab:audit}. The bound covers the realized residual harm on 7
of 9 targets (78\,\%), below the 90\,\% we required, so the
pre-registered claim is refuted. The two misses are nuScenes, a real
shift, and Palo Alto, where 8.0\,\% against 7.9\,\% on 253 scenarios is within
sampling noise. The audit is valid when the target resembles the audit
split and must be recomputed on target-like labeled data otherwise.

\begin{table}[h]
\caption{Realized residual harm $O$ and total harm $O_{\mathrm{tot}}$
against the audit upper bounds (7.9\,\% and 9.4\,\%). Rates in \%.}
\label{tab:audit}
\centering
\scriptsize\setlength{\tabcolsep}{2.5pt}
\begin{tabular}{lccc}
\toprule
Target & $O$ & $O_{\mathrm{tot}}$ & Covered \\
\midrule
Fresh test (in distribution) & 4.8 & 7.2 & yes \\
Austin & 7.0 & 8.6 & yes \\
Dearborn & 4.4 & 5.1 & yes \\
Miami & 6.8 & 8.3 & yes \\
Palo Alto & 8.0 & 10.0 & no \\
Pittsburgh & 6.3 & 8.4 & yes \\
Washington, DC & 6.7 & 8.1 & yes \\
nuScenes & 15.6 & 16.4 & no \\
Waymo & 4.3 & 8.3 & yes \\
\bottomrule
\end{tabular}
\end{table}

\section{Corrected Validity and the Hypothesis Ledger}
\label{sec:s-ledger}
The violation frequency reported in the main paper counts resplits whose
test-half miss exceeds $\alpha$, which also counts test-sample noise. We
also report a corrected count that only counts resplits whose one-sided
95\,\% Wilson lower bound exceeds $\alpha$, and call a method \emph{valid} if
the uncorrected count is at most $\delta$, \emph{indeterminate} if only the
corrected count is, and \emph{violated} otherwise (Amendment~11).
Table~\ref{tab:valid} gives every method. The separation of Learn-then-Test
from the uncertified methods survives the correction, whereas the
near-boundary cases (latency-aware calibration, group-wise resplits)
become indeterminate rather than violated.

\begin{table}[h]
\caption{Violation frequency over 200 resplits, uncorrected / corrected,
and verdict. Rates in \%.}
\label{tab:valid}
\centering
\scriptsize\setlength{\tabcolsep}{2.5pt}
\begin{tabular}{lcc}
\toprule
Method & Unc. / corr. & Verdict \\
\midrule
LTT (geometric) & 9.5 / 0.0 & valid \\
Tuned geometric & 43.5 / 15.0 & violated \\
CRC & 41.0 / 12.0 & violated \\
Tuned T1 confidence & 50.5 / 14.0 & violated \\
LTT latency-aware, 0.5\,s & 12.0 / 1.5 & indeterminate \\
LTT latency-aware, 1.0\,s & 11.5 / 0.5 & indeterminate \\
Zero-latency LTT at 0.5\,s & 92.0 / 53.0 & violated \\
Zero-latency LTT at 1.0\,s & 100 / 100 & violated \\
LTT, fresh data & 7.5 / 0.5 & valid \\
LTT, second predictor & 9.0 / 0.0 & valid \\
LTT, Wayformer & 7.0 / 0.0 & valid \\
LTT, spatial groups (150\,m) & 17.0 / 1.0 & indeterminate \\
TTC, LTT & 1.0 / 1.0 & valid \\
CDT (100 resplits) & 59.0 / 0.0 & indeterminate \\
ACI (100 resplits) & 0.0 / 0.0 & valid \\
\bottomrule
\end{tabular}
\end{table}

Every claim of Amendment~11 and its addendum was fixed before it was run, and the ledger records each verdict.

\begin{table}[h]
\caption{Amendment~11 claims and verdicts.}
\label{tab:ledger}
\centering
\scriptsize\setlength{\tabcolsep}{2.5pt}
\begin{tabular}{p{5.2cm}c}
\toprule
Claim & Verdict \\
\midrule
P1: the predictor-based trigger stops less than a physics TTC trigger & refuted \\
N-A: prevalence-corrected certificate meets the target on at least 7 of 8 shifts and beats plain LTT & refuted \\
N-B: latency-aware LTT is not violated at 0.5 and 1.0\,s & holds \\
N-C: the outcome audit bound covers at least 90\,\% of targets & refuted \\
CSWC (i): outcome certificate valid at 7.5\,\% & holds \\
CSWC (ii): learned trigger stops less than geometric and TTC at zero latency & refuted \\
CSWC (iii): prevalence correction helps the learned trigger under shift & refuted \\
11b: learned trigger valid and cheapest at 0.5 and 1.0\,s & holds \\
11c M1: the predictor adds lead time beyond physics features & refuted \\
11c R1: the learned trigger replicates at 0.5\,s on a second AutoBot and Wayformer & holds \\
\bottomrule
\end{tabular}
\end{table}

\section{Mechanism, Replication and Shift for the Counterfactual Trigger}
\label{sec:s-mech}
These analyses were pre-registered in Amendment~11c.

\emph{Mechanism (claim refuted).} To test whether the predictor adds lead
time that physics features lack, we trained three learned triggers with
delay-shifted labels at a latency of 0.5\,s (one tick): all features, physics
features only (time-to-collision, headway distance, agent count, tick index
and scenario covariates) and predictor features only (geometric and
confidence scores, predicted gap, tick index and covariates). All three
certify the 7.5\,\% target. Table~\ref{tab:mech} shows that the full trigger
does not stop fewer scenarios than the physics-only trigger (paired CI of the
difference $[-0.1, 1.6]$ points), and the predictor-only trigger stops fewer
(4.0 points, CI $[2.7, 5.2]$, though it sits closer to the target). The
pre-registered claim is therefore refuted: the gain over a plain TTC
threshold comes from learning on counterfactual, delay-shifted labels, and
either feature group suffices.

\begin{table}[h]
\caption{Feature-group ablation of the learned trigger at 0.5\,s of latency
(calibrate 904, test 2{,}900). $O$: realized residual harm; $O_{\mathrm{tot}}$:
with induced collisions. Rates in \%.}
\label{tab:mech}
\centering
\scriptsize\setlength{\tabcolsep}{3pt}
\begin{tabular}{lccc}
\toprule
Features & $O$ & $O_{\mathrm{tot}}$ & Unnec. stops \\
\midrule
All & 4.9 & 6.5 & 25.5 \\
Physics only & 4.9 & 6.6 & 24.8 \\
Predictor only & 5.7 & 7.3 & 21.5 \\
\bottomrule
\end{tabular}
\end{table}

\emph{Replication (claim holds).} With the second AutoBot and with Wayformer
the calibration arm is split into 452 training and 452 calibration scenarios
and the 904 test scenarios are held out. At 0.5\,s the learned trigger meets
the target (3.3\,\% and 4.3\,\%) and stops 49.3\,\% and 43.0\,\% against 84.8\,\%
for both the geometric and the TTC score (paired CIs exclude 0). At zero
latency it is worse than the geometric score on the second AutoBot (46.2\,\%
against 39.5\,\%, CI $[3.1, 10.4]$) and better on Wayformer (44.8\,\% against
57.0\,\%); both beat TTC, which only certifies at always firing with 452
calibration scenarios. The stop rates are higher than in the main split
because the training and calibration sets are half as large.

\emph{Shift (descriptive).} Table~\ref{tab:scoreshift} gives the realized
residual harm of the outcome-certified threshold of each score on the eight
shift targets of the main paper (zero latency, 7.5\,\% target). Each score
meets the target on 6 of 8. The learned trigger fails Miami and nuScenes, the
geometric score nuScenes and Washington, DC, and TTC Miami and nuScenes. No
score survives shift better than the others, and TTC is usually cheapest.

\begin{table}[h]
\caption{Realized residual harm $O$ and unnecessary stops of the
outcome-certified threshold on each shift target (zero latency). Rates in \%.}
\label{tab:scoreshift}
\centering
\scriptsize\setlength{\tabcolsep}{2.5pt}
\begin{tabular}{lcccccc}
\toprule
 & \multicolumn{2}{c}{Learned} & \multicolumn{2}{c}{Geometric} & \multicolumn{2}{c}{TTC} \\
\cmidrule(lr){2-3}\cmidrule(lr){4-5}\cmidrule(lr){6-7}
Target & $O$ & Stops & $O$ & Stops & $O$ & Stops \\
\midrule
Austin & 6.3 & 8.3 & 6.6 & 26.7 & 4.8 & 8.6 \\
Dearborn & 2.7 & 7.5 & 3.5 & 18.3 & 2.9 & 5.3 \\
Miami & 10.8 & 7.0 & 7.4 & 35.1 & 9.2 & 9.4 \\
Palo Alto & 5.1 & 8.9 & 5.1 & 22.9 & 4.5 & 6.8 \\
Pittsburgh & 5.8 & 13.0 & 5.5 & 37.9 & 6.1 & 13.8 \\
Washington, DC & 7.4 & 7.4 & 7.6 & 27.1 & 6.7 & 8.6 \\
nuScenes & 8.2 & 59.8 & 12.3 & 34.8 & 8.8 & 20.7 \\
Waymo & 4.3 & 30.4 & 3.8 & 48.1 & 4.8 & 9.9 \\
\bottomrule
\end{tabular}
\end{table}

\section{Outcome Partition of the Fallback}
\label{sec:s-partition}
The table partitions every scenario by what happens after the fallback fires, as described in the main paper.

\begin{table}[h]
\caption{Outcome partition (\% of scenarios). Residual = miss + too late +
new harm. Induced: not-at-fault collision the no-fallback run avoids.}
\label{tab:cost}
\centering
\scriptsize
\setlength{\tabcolsep}{3pt}
\begin{tabular}{llccccc}
\toprule
Data & Trigger & Miss & Too late & Prev. & New harm & Resid. / Induced \\
\midrule
AV2 & No fallback & 15.1 & 0.0 & 0.0 & 0.0 & 15.1 / 0.0 \\
AV2 & LTT (certified) & 3.7 & 1.7 & 9.8 & 1.1 & 6.5 / 2.6 \\
AV2 & Tuned geometric & 4.3 & 1.7 & 9.1 & 0.9 & 6.9 / 2.3 \\
AV2 & Always fire & 0.0 & 1.0 & 14.1 & 3.1 & 4.1 / 6.5 \\
nuScenes & LTT (AV2-cert.) & 13.2 & 1.0 & 8.2 & 1.7 & 15.6 / 2.4 \\
nuScenes & Always fire & 5.0 & 1.6 & 16.2 & 3.8 & 9.7 / 6.4 \\
\bottomrule
\end{tabular}
\end{table}

\end{document}
